\section{Results}\label{sec:results}

This section reports the 2017 held-out year, and every number comes from the frozen ledger. Errors follow $e=B-L$, so a negative bias means that a baseline under-predicts the load and underpays the participant. With $S=\,${{sample.n_sites|int}} sites the smallest attainable one-sided exact p-value is $2^{-{{sample.n_sites|int}}}$, and a Holm-adjusted value cannot fall below five times that floor. Section~\ref{sec:results:accuracy} answers RQ1 (H1 to H3), Section~\ref{sec:results:h4} answers RQ2 (H4) and Section~\ref{sec:results:h5} answers RQ3 (H5).

\subsection{Sample}\label{sec:results:sample}

The final sample has {{sample.n_buildings|int}} buildings at {{sample.n_sites|int}} sites (Fig.~\ref{fig:sample}; eligibility by site in Supplement Table~\ref{tabS:data}). The hot-day rule R1 yields {{sample.n_events.R1|int}} building-events and the random-weekday rule R3 yields {{sample.n_events.R3|int}}. Every rule has a baseline for at least {{sample.support_share.ny0|pct1}} of the R1 events. CG-adj is missing for {{derived.cgA.n_events_missing_eagle|int}} Eagle events that lack five complete donors, so H3 rests on {{H3.logratio.cgA_pjmA:n_events|int}} events. Lamb keeps {{sample.excl.Lamb.n_after|int}} buildings after the flat-meter filter, and only {{sample.event_days.Lamb|int}} of its {{sample.event_days_designated.Lamb|int}} designated event days carry scored building-events, in {{sample.heat_episodes.Lamb|int}} heat episodes, so its site values rest on few buildings and days. Supplement Table~\ref{tabS:events} gives event days and episodes by site.

\subsection{Hot-day accuracy and the tests H1 to H3}\label{sec:results:accuracy}

Adjusted rules were more accurate on hot days than unadjusted ones; this summary describes the named baselines and their metrics and is not a tested attribution of error to information set or rule origin (Table~\ref{tab:accuracy}, Fig.~\ref{fig:accuracy}). The building-pooled median relMAE of PJM-adj was {{method.pjmA.R1.relMAE_median|pct1}} against {{method.pjm0.R1.relMAE_median|pct1}} for unadjusted PJM; its site median was {{derived.T4.pjmA.relMAE|pct1}}, with a 95\,\% order-statistic interval from {{derived.T4.pjmA.relMAE:ci_low|pct1}} to {{derived.T4.pjmA.relMAE:ci_high|pct1}}. The other adjusted rules, except TOWT-70-adj ({{method.towt70A.R1.relMAE_median|pct1}}), ranged from {{method.isoneA.R1.relMAE_median|pct1}} (ISO-NE-adj) to {{method.nyA.R1.relMAE_median|pct1}} (NYISO-adj), and unadjusted day-matching rules from {{method.pjm0.R1.relMAE_median|pct1}} (PJM) to {{method.mid8.R1.relMAE_median|pct1}} (Mid8). Fixed-parameter LightGBM-morn and LightGBM-donor reached {{method.lgbM.R1.relMAE_median|pct1}} and {{method.lgbD.R1.relMAE_median|pct1}}, and an essentially unchanged {{method.lgbM_hpo.R1.relMAE_median|pct1}} and {{method.lgbD_hpo.R1.relMAE_median|pct1}} after the pre-registered search ({{hpo.lgbM.trials|int}} trials per variant, tuned on 2016 only). The tuned MLP variants (pre-registered {{hpo.mlpM.trials|int}}-trial search on 2016, {{method.mlpM.R1.n_seeds|int}} training seeds) reached {{method.mlpM.R1.relMAE_median|pct1}} (MLP-morn) and {{method.mlpD.R1.relMAE_median|pct1}} (MLP-donor), about the accuracy of the tuned LightGBM variants with the same information set. PJM-adj stayed more accurate than all four tuned learned baselines.

Bias told the same story. Unadjusted PJM had a pooled median bias of {{method.pjm0.R1.bias_median|pct1}} on hot days and PJM-adj {{method.pjmA.R1.bias_median|pct1}}. In Granderson's sign convention these are a normalized mean bias error of {{method.pjm0.R1.nmbe_granderson|pct1}} and {{method.pjmA.R1.nmbe_granderson|pct1}}, so the unadjusted rule under-predicted in the building-pooled summary, in line with prior peak-window results \cite{granderson2021}. Its site-median hot-day bias was {{derived.T4.pjm0.bias|pct1}} (interval {{derived.T4.pjm0.bias:ci_low|pct1}} to {{derived.T4.pjm0.bias:ci_high|pct1}}), so the under-prediction is not a site-level result. The previous-weekday baseline had a relMAE of {{method.prev.R1.relMAE_median|pct1}} in the pooled summary but {{derived.T4.prev.relMAE|pct1}} against {{derived.T4.pjm0.relMAE|pct1}} for PJM in the site-median summary. The two orders disagree and the intervals overlap, so no ranking of Prev and PJM is claimed. Over building-events, unadjusted PJM underpaid {{derived.X4.pjm0.share_underpaid|pct1}} of hot-day cases, Prev {{derived.X4.prev.share_underpaid|pct1}} and PJM-adj {{derived.X4.pjmA.share_underpaid|pct1}}. These shares are descriptive, and no rule is declared best.

\paragraph{H1 (replication)}The hot-day bias of unadjusted PJM was more negative than its random-weekday bias. The Hodges-Lehmann site-level difference was {{H1.delta_bias.pjm0|.3f}} (95\,\% interval {{H1.delta_bias.pjm0:ci_low|.3f}} to {{H1.delta_bias.pjm0:ci_high|.3f}}; exact one-sided p = {{H1.p_signedrank|.5f}}; Holm p = {{H1.p_holm|.5f}}), and {{derived.H1.n_sites_negative|int}} of {{sample.n_sites|int}} sites had a negative difference (Supplement Fig.~\ref{figS:h1bias}). Lamb was the exception. Leaving out any one site moved the estimate between {{H1.loso_min|.3f}} and {{H1.loso_max|.3f}}. The pooled bias on random weekdays (R3) was {{method.pjm0.R3.bias_median|pct1}}, and the secondary NYISO and CAISO or ISO-NE contrasts pointed the same way ({{H1s.delta_bias.ny|.3f}} and {{H1s.delta_bias.caiso|.3f}}), each with a Holm p-value of {{H1s.p_holm_family.ny|.5f}} within the pre-specified secondary rule family.

\paragraph{H2 (replication)}The day-of adjustment reduced the hot-day error of PJM, which replicates the finding of Granderson et al.\ \cite{granderson2021}. The log ratio of relMAE, PJM-adj over PJM, was {{H2.logratio.pjmA_pjm0|.3f}} (95\,\% interval {{H2.logratio.pjmA_pjm0:ci_low|.3f}} to {{H2.logratio.pjmA_pjm0:ci_high|.3f}}), negative at {{derived.H2.n_sites_negative|int}} of {{sample.n_sites|int}} sites. The one-sided p-value, {{H2.p_signedrank|.6f}}, sits at the $2^{-S}$ floor, so the test shows the sign of the effect, not its size, and the Holm p-value is {{H2.p_holm|.5f}}. Eagle gained almost nothing ({{H2.site.Eagle|.3f}}). Removing single sites moved the estimate between {{H2.loso_min|.3f}} and {{H2.loso_max|.3f}} (Supplement Fig.~\ref{figS:h2}). The day-of adjustment also reduced the error of the NYISO, CAISO and ISO-NE rules, each with a within-family Holm p-value of {{H2s.p_holm_family.ny|.5f}}.

\paragraph{H3 (equivalence of the site location)}Under the signed-rank location (symmetry) assumption, CG-adj and PJM-adj were equivalent in the location of the site log error ratios, within the pre-set ten percent margin ($\delta={{H3.margin|.4f}}$ on the log scale), over {{H3.logratio.cgA_pjmA:n_events|int}} building-events without the {{derived.cgA.n_events_missing_eagle|int}} Eagle events that lack donors. The estimate was {{H3.logratio.cgA_pjmA|.3f}} with a 90\,\% interval from {{H3.logratio.cgA_pjmA:ci_low|.3f}} to {{H3.logratio.cgA_pjmA:ci_high|.3f}}, the TOST p-value was {{H3.p_tost|.4f}}, and the decision was \emph{{{H3.decision|raw}}}. The Holm-adjusted value, {{H3.p_holm|.4f}}, equals the raw one because H3 is last in the Holm order. {{derived.H3.n_sites_outside_margin|int}} of {{sample.n_sites|int}} sites lay outside the margin (Fig.~\ref{fig:h2h3}), so equivalence does not hold site by site, and the share inside was {{derived.H3.share_sites_inside_margin|pct1}} (exact binomial 95\,\% interval {{derived.H3.share_sites_inside_margin:ci_low|pct1}} to {{derived.H3.share_sites_inside_margin:ci_high|pct1}}). The conditional bootstrap interval, {{H3.ci_boot_conditional:ci_low|.3f}} to {{H3.ci_boot_conditional:ci_high|.3f}}, crosses the lower margin. CG-adj is PJM-adj rescaled by the donors' ratio of metered load to their own PJM-adj baselines (Eq.~\eqref{eq:cga}), and its donors share the target site's day-to-day shocks. The similarity is therefore partly by construction, and the test cannot show agreement where donors and target face different shocks. The share inside, descriptive only, was {{derived.H3.share_inside_margin.m05|pct1}} at a five percent margin and {{derived.H3.share_inside_margin.m15|pct1}} at 15 percent. The two rules are not declared interchangeable for individual buildings or sites.

\subsection{Pre-event shaping (H4)}\label{sec:results:h4}

In the four-hour injected stress test, PJM-adj lost its accuracy advantage over SCM-ridge at a break-even pre-event increase of {{H4.median_site_pcross|pct1}} at the median site (order-statistic interval {{H4.median_site_pcross:ci_low|pct1}} to {{H4.median_site_pcross:ci_high|pct1}}; coverage at least 95\,\%, exact {{derived.H4.ci_exact_coverage|pct1}}). This statistic is the median of the {{sample.n_sites|int}} site break-evens. The increase is injected, not observed, and the realistic band is an assumption converted symbolically from field energy penalties. The interval starts at zero because {{derived.H4.n_sites_pcross_zero|int}} sites, Eagle and Lamb, have a break-even of zero: SCM-ridge was already at least as accurate there without any increase. All {{derived.H4.n_sites_pcross_lt_threshold|int}} of {{sample.n_sites|int}} sites lay below the pre-set threshold of 0.10, and the largest site value was {{derived.H4.max_site_pcross|pct1}} (Fig.~\ref{fig:precool}a). The exact sign test gave {{H4.p_sign|.6f}}, the one-sided floor, and the Holm p-value was {{H4.p_holm|.5f}}. Leaving out one site moved the median between {{H4.loso_min|pct1}} and {{H4.loso_max|pct1}}. The building-pooled four-hour break-even, a separate and secondary estimand, was {{H4.pcross.pjmA_scmR.4h|pct1}} (interval {{H4.pcross.pjmA_scmR.4h:ci_low|pct1}} to {{H4.pcross.pjmA_scmR.4h:ci_high|pct1}}).

For the two-hour shape the conclusion does not transfer. Its result is building-pooled: the break-even was {{pcross.pjmA.2h|pct1}} with an interval of {{pcross.pjmA.2h:ci_low|pct1}} to {{pcross.pjmA.2h:ci_high|pct1}}, and the interval includes the threshold of 0.10. No site-level two-hour test was specified or run. At an injected increase of 0.10 the pooled relMAE of PJM-adj rose by {{incr_sens.4h.p0.1|pct1}} of window load for the four-hour shape and by {{incr_sens.2h.p0.1|pct1}} for the two-hour shape, a ratio of {{derived.incr_sens.ratio_2h_4h|.2f}}. That ratio is below the overlap share of one third, so the first-order expression is a guide to the direction and not a fitted prediction. SCM-ridge is immune by construction to the specified four-hour injection, which is a design property and not a finding, and its error did not change with $p$ (Fig.~\ref{fig:precool}).

The primary shape raises load in the four hours before the event and so covers all three PJM adjustment hours, which makes the H4 break-even a worst-case placement bound for PJM-adj. When the shaping is shifted earlier so that it overlaps only the first of the three adjustment hours (exploratory), the median site break-even, taken as in H4 against SCM-ridge without an increase, rises to {{exploratory.early_shape.a.pcross_frozenstyle.median_site|pct1}}, about three times the primary value. Shaping that touches no adjustment hour leaves PJM-adj unchanged by construction (Supplement Fig.~\ref{figS:placement}). In practitioner units, and as a property of the injected test only, the extra pre-event energy at the break-even equals {{derived.H4.energy_offset_share.median_site|pct1}} of the energy shed by a ten percent curtailment at the median site. At that increase PJM-adj would pay about {{derived.H4.overpay_multiplier.median_site|.2f}} times the true shed at the median site, assuming that the curtailment leaves the adjustment hours unchanged and that payment is proportional to the baseline minus metered load.

The 10\,\%-error thresholds $p^{\star}$ of the other rules (Fig.~\ref{fig:pstar}) follow the overlap of hours read and hours shaped. Under the four-hour shape the pooled relMAE of the adjusted rules reached the curtailment value at an increase between {{pstar.wmA.4h|pct1}} (WM-adj) and {{pstar.caisoA.4h|pct1}} (CAISO-adj), with PJM-adj at {{pstar.pjmA.4h|pct1}} (interval {{pstar.pjmA.4h:ci_low|pct1}} to {{pstar.pjmA.4h:ci_high|pct1}}). NYISO-adj crossed at {{pcross.nyA.4h|pct1}}, because its error with no injection ({{method.nyA.R1.relMAE_median|pct1}}) is already close to that of SCM-ridge ({{method.scmR.R1.relMAE_median|pct1}}). For the two-hour shape only PJM-adj ({{pstar.pjmA.2h|pct1}}), CG-adj ({{pstar.cgA.2h|pct1}}) and ISO-NE-adj ({{pstar.isoneA.2h|pct1}}) reached the curtailment value within the grid; NYISO-adj stays flat because its adjustment hours fall outside the shaped hours.

Two further checks follow. In the pre-specified donor-contamination stress test, a share $q$ of the donors also curtailed and shaped alongside the participants, with an increase of 0.05 injected. SCM-ridge's relMAE advantage over PJM-adj, which reads no donors, was {{derived.contam.edge_scmR_vs_pjmA.q0|pct1}} of window load with clean donors and {{derived.contam.edge_scmR_vs_pjmA.q0.25|pct1}} with a quarter of them participating, and it turned to {{derived.contam.edge_scmR_vs_pjmA.q0.5|pct1}} with half. Linear interpolation on this four-point grid puts the crossing near $q={{derived.contam.q_cross_scmR_pjmA|.2f}}$, a descriptive value. The H4 break-even therefore describes uncontaminated donors, the case that favors SCM-ridge. The one post-freeze exploratory correlation between the site break-even and the site gap in relMAE at $p=0$ was {{exploratory.x6.rho|.2f}}, with an unadjusted exact permutation p-value of {{exploratory.x6.p_perm|.3f}} over {{exploratory.x6.n_sites|int}} sites. The sign is the expected one, but the correlation is not distinguishable from zero.

\subsection{Repeated dispatch (H5)}\label{sec:results:h5}

In the injected repeated-dispatch stress test, PJM-adj had a smaller site-median absolute error than SCM-ridge at every tested dispatch count in both event windows (Table~\ref{tab:h5windows}, Fig.~\ref{fig:dispatch}a). The curtailment is injected, and the metered series is the no-curtailment counterfactual by assumption only. In the confirmatory 14:00 to 18:00 window, the site median of the absolute deviation of the estimated curtailment ratio was {{derived.H5.absdev_sitemedian.scmR.K10|pct1}} for SCM-ridge against {{derived.H5.absdev_sitemedian.pjmA.K10|pct1}} for PJM-adj at $K=10$. At $K=40$ the values were {{derived.H5.absdev_sitemedian.scmR.K40|pct1}} and {{derived.H5.absdev_sitemedian.pjmA.K40|pct1}}. PJM-adj was lower at {{derived.H5.n_sites_pjmA_lower.K10.w1418|int}} of {{sample.n_sites|int}} sites at every $K$. In the pre-specified secondary 16:00 to 20:00 window the values were {{derived.H5.absdev_sitemedian_w1620.scmR.K10|pct1}} against {{derived.H5.absdev_sitemedian_w1620.pjmA.K10|pct1}} at $K=10$ and {{derived.H5.absdev_sitemedian_w1620.scmR.K40|pct1}} against {{derived.H5.absdev_sitemedian_w1620.pjmA.K40|pct1}} at $K=40$, with PJM-adj lower at {{derived.H5.n_sites_pjmA_lower.K10.w1620|int}} to {{derived.H5.n_sites_pjmA_lower.K40.w1620|int}} sites. The building-pooled summary agreed in both windows, at {{derived.H5.levels_pjmA_lower_all_K_pooled.w1418|int}} and {{derived.H5.levels_pjmA_lower_all_K_pooled.w1620|int}} of 4 tested counts. CG-adj sat close to PJM-adj, and unadjusted PJM far above both. $K=10$ is near the upper end of documented annual event counts \cite{kiliccote2010,nyiso2024dr}, and larger $K$ is a stress test.

In the 14:00 to 18:00 window, the confirmatory slope contrast agreed with the levels. The Hodges-Lehmann difference in site slopes, PJM-adj minus SCM-ridge, was {{H5.eta_diff|.4f}} (95\,\% interval {{H5.eta_diff:ci_low|.4f}} to {{H5.eta_diff:ci_high|.4f}}), where a negative value means that SCM-ridge's slope is larger. The two-sided exact p-value was {{H5.p_signedrank|.5f}} and the Holm p-value {{H5.p_holm|.5f}}; the raw value is not at its two-sided floor. SCM-ridge's slope was larger at {{derived.H5.n_sites_scmR_faster|int}} of {{sample.n_sites|int}} sites (Fig.~\ref{fig:dispatch}c), and removing any single site left the estimate between {{H5.loso_min|.4f}} and {{H5.loso_max|.4f}}. The median of the site slopes was {{H5.eta_abs.median_of_site_slopes.scmR|.4f}} per dispatch for SCM-ridge and {{H5.eta_abs.median_of_site_slopes.pjmA|.4f}} for PJM-adj, so PJM-adj's absolute error also grows. The test is two-sided, and the absolute definition was fixed on the development year (Section~\ref{sec:data:plan}).

In the pre-specified secondary 16:00 to 20:00 window, the slope contrast was inconclusive. The difference was {{exploratory.h5_win1620.eta_diff.pjmA_minus_scmR|.4f}} (interval {{exploratory.h5_win1620.eta_diff.pjmA_minus_scmR:ci_low|.4f}} to {{exploratory.h5_win1620.eta_diff.pjmA_minus_scmR:ci_high|.4f}}; unadjusted two-sided p = {{exploratory.h5_win1620.p_signedrank_two_sided_unadjusted|p}}, outside the Holm family), in the same direction but with an interval that includes zero (Table~\ref{tab:h5windows}). The relative growth from $K=10$ to $K=40$ also reversed there: building-pooled, PJM-adj's absolute error grew by {{derived.H5.relgrowth_K10_K40_pooled.pjmA.w1620|pct0}} and SCM-ridge's by {{derived.H5.relgrowth_K10_K40_pooled.scmR.w1620|pct0}} ({{derived.H5.relgrowth_K10_K40_pooled.pjmA.w1418|pct0}} and {{derived.H5.relgrowth_K10_K40_pooled.scmR.w1418|pct0}} at 14:00 to 18:00), and the site medians grew about equally ({{derived.H5.relgrowth_K10_K40_sitemedian.pjmA.w1620|pct0}} and {{derived.H5.relgrowth_K10_K40_sitemedian.scmR.w1620|pct0}}). In the 2016 development year, the 16:00 to 20:00 contrast had the same sign and an interval that excluded zero (development, not a result). The level ordering therefore holds in both windows, whereas the confirmatory slope finding is stated for the 14:00 to 18:00 window only. An exploratory paired contrast of the two windows is in Supplement Section~\ref{sec:S1:inventory}.

The signed curves show a different pattern (Fig.~\ref{fig:dispatch}b). Unadjusted PJM eroded in the signed sense: its median ratio of estimated to injected curtailment fell from {{H5.rho.pjm0.K10|.2f}} at $K=10$ to {{H5.rho.pjm0.K40|.2f}} at $K=40$, the latter a stress level. PJM-adj stayed near one ({{H5.rho.pjmA.K10|.2f}} to {{H5.rho.pjmA.K40|.2f}}), and SCM-ridge stayed above one, near {{H5.rho.scmR.K10|.2f}}, throughout. The secondary signed-error contrast was not distinguishable from zero ({{H5s.eta_signed_diff.pjmA_minus_scmR|.4f}}, interval {{H5s.eta_signed_diff.pjmA_minus_scmR:ci_low|.4f}} to {{H5s.eta_signed_diff.pjmA_minus_scmR:ci_high|.4f}}, p = {{H5s.eta_signed_diff.pjmA_minus_scmR:p_raw|.3f}}), so SCM-ridge's error grows in size and not in a consistent direction. Its pre-period was older than the PJM lookback at every $K$ ({{H5.age_days.scmR_preperiod.K40|.1f}} against {{H5.age_days.pjm_lookback.K40|.1f}} days at $K=40$; Supplement Fig.~\ref{figS:age}). This diagnostic is not proof of a mechanism, but it is consistent in sign with the dispersion growth under shifted lookbacks that Granderson et al.\ report \cite{granderson2023}.

PJM had no fallbacks through $K=30$, and they reached {{H5.fallback.PJM.K40|pct1}} of building-events at $K=40$, while CAISO had fewer than ten eligible days in {{H5.fallback.CAISO_lt10.K40|pct1}} of them (Supplement Table~\ref{tabS:fallback}). When all heat-wave days were dispatched as well (median {{heatwave.n_dispatched_median|int}} days), the median ratios were {{heatwave.pjmA.rho_median|.2f}} for PJM-adj and {{heatwave.scmR.rho_median|.2f}} for SCM-ridge, with relMAE of {{heatwave.pjmA.relMAE_median|pct1}} and {{heatwave.scmR.relMAE_median|pct1}}. No slope contrast was computed on that subset.

\subsection{Payment consequences}\label{sec:results:payment}

The payment error ratio is the error in units of the value of a ten percent curtailment. PJM-adj had a median error of about {{method.pjmA.R1.PER_c10_median|.2f}} times that value and unadjusted PJM about {{method.pjm0.R1.PER_c10_median|.2f}} times (Fig.~\ref{fig:payment}). The error exceeded half of the value in {{method.pjmA.R1.share_PER_gt_0_5|pct1}} of PJM-adj events and in {{method.pjm0.R1.share_PER_gt_0_5|pct1}} of unadjusted PJM events. SCM-ridge had {{method.scmR.R1.PER_c10_median|.2f}} and {{method.scmR.R1.share_PER_gt_0_5|pct1}}, and CG-adj {{method.cgA.R1.PER_c10_median|.2f}} and {{method.cgA.R1.share_PER_gt_0_5|pct1}}. Even the most accurate baselines carry a median hot-day payment error of roughly half the value of a ten percent curtailment (other fractions are in the Supplement).

\subsection{Sensitivity of the confirmatory tests}\label{sec:results:robust}

Table~\ref{tab:confirmatory} summarizes the five tests (sensitivity analyses in Supplement Fig.~\ref{figS:robustness}). All five confirmatory decisions met their Holm-adjusted criteria: H1, H2, H4 and H5 rejected their null hypotheses, and H3 met its equivalence criterion. The wild cluster bootstrap targets the site mean, a different estimand, and gave p-values of {{H1.p_wcb|.4f}} (H1), {{H2.p_wcb|.4f}} (H2), {{H4.p_wcb|.4f}} (H4) and {{H5.p_wcb|.4f}} (H5); H3 has none. These values can fall below $2^{-S}$ because the bootstrap draws one of six Webb weights per site for the building-level residuals. It thus draws from $6^{S}$ rather than $2^{S}$ possible weight vectors, and it adds no independent sites. Claims are restricted to these sites.

The order of the method levels holds in the seven-site R2-ISO sample, where PJM-adj, CG-adj and SCM-ridge had relMAE of {{method.pjmA.R2ISO.relMAE_median|pct1}}, {{method.cgA.R2ISO.relMAE_median|pct1}} and {{method.scmR.R2ISO.relMAE_median|pct1}}. It also holds in the 16:00 to 20:00 window, where they had {{method.pjmA.w1620.relMAE_median|pct1}}, {{method.cgA.w1620.relMAE_median|pct1}} and {{method.scmR.w1620.relMAE_median|pct1}} (Fig.~\ref{figS:sensitivity}). No break-even was computed for R2-ISO. These and the remaining reruns are pre-specified secondary analyses. In the 16:00 to 20:00 window, the H4 break-even stayed low at the median site ({{exploratory.h4_win1620.median_site_pcross|pct1}}, interval {{exploratory.h4_win1620.median_site_pcross:ci_low|pct1}} to {{exploratory.h4_win1620.median_site_pcross:ci_high|pct1}}; unadjusted exact sign p = {{exploratory.h4_win1620.p_sign_unadjusted|.5f}}, at the floor), and its H5 contrast is reported in Section~\ref{sec:results:h5}. With the opposite daylight-saving mapping, H4 and H5 were essentially unchanged (median break-even {{exploratory.dst_alt.h4.median_site_pcross|pct1}}, slope difference {{exploratory.dst_alt.h5.eta_diff.pjmA_minus_scmR|.4f}}, p = {{exploratory.dst_alt.h5.p_signedrank_two_sided_unadjusted|p}}). Applying the flat-meter filter to non-event days kept {{nonevent.sample.n_buildings_R1|int}} buildings, and the H1, H2 and H3 estimates were {{nonevent.H1.est|.3f}}, {{nonevent.H2.est|.3f}} and {{nonevent.H3.est|.3f}}. Switching the filter off gave {{flatoff.H1.est|.3f}}, {{flatoff.H2.est|.3f}} and {{flatoff.H3.est|.3f}}, with H1 then over {{flatoff.H1.est:n_sites|int}} sites. The opposite daylight-saving mapping also left the PJM-adj relMAE at {{dst_alt.pjmA.R1.relMAE_median|pct1}} and the H2 estimate at {{dst_alt.H2.logratio.pjmA_pjm0|.3f}}. Adding the rebound term left the unadjusted baselines unchanged by construction.

Ablations fall into two tiers. The adjustment itself is the strong choice: turning it off raised the PJM-adj relMAE from {{ablation.adjustment.pjmA.on.relMAE_median|pct1}} to {{ablation.adjustment.pjmA.off.relMAE_median|pct1}}. Other adjustment hours and a cap moved it by less, to values between {{ablation.adjhours.pjmA.s-3..s-2.relMAE_median|pct1}} and {{ablation.adjhours.pjmA.s-4..s-3.relMAE_median|pct1}} (cap {{ablation.cap.pjmA.on.relMAE_median|pct1}}; Supplement Table~\ref{tabS:secondary}). Cutting the CG-adj donor pool to five buildings gave a log error ratio of {{ablation.donor.k5.logratio_vs_wholepool|.3f}} against the whole pool (interval {{ablation.donor.k5.logratio_vs_wholepool:ci_low|.3f}} to {{ablation.donor.k5.logratio_vs_wholepool:ci_high|.3f}}).

\subsection{Exploratory runs}\label{sec:results:explore}

These runs were decided after the 2017 confirmatory results were seen, as the post-freeze register records (Supplement Section~\ref{sec:S1:dev}). They stay outside the Holm family, and every p-value below is unadjusted. SCM-ridge with the longer pre-period, $W=10$ days chosen on 2016 accuracy, gave a slope difference of {{exploratory.x7.h5_W10.eta_diff.pjmA_minus_scmR|.4f}} (interval {{exploratory.x7.h5_W10.eta_diff.pjmA_minus_scmR:ci_low|.4f}} to {{exploratory.x7.h5_W10.eta_diff.pjmA_minus_scmR:ci_high|.4f}}; p = {{exploratory.x7.h5_W10.p_signedrank_two_sided_unadjusted|p}}), so the H5 direction in the 14:00 to 18:00 window holds for that window length.

Two further runs probe where PJM-adj is exposed. In early-shaping variant a, the shaping was shifted earlier so that it overlaps only the first of the three adjustment hours. Its median site break-even, again against SCM-ridge without an increase as in H4 (Section~\ref{sec:results:h4}; interval {{exploratory.early_shape.a.pcross_frozenstyle.median_site:ci_low|pct1}} to {{exploratory.early_shape.a.pcross_frozenstyle.median_site:ci_high|pct1}}) was not below the pre-set threshold of 0.10 at most sites (exact sign p = {{exploratory.early_shape.a.pcross_frozenstyle.p_sign_unadjusted|p}}). PJM-adj is thus most exposed when the shaping falls in the hours its adjustment reads. Variant b, in which no adjustment hour is shaped, leaves PJM-adj unchanged by construction and was not run. In the mirror test, a pre-event decrease of the same size, the median site break-even magnitude was {{exploratory.neg_p.median_site_pcross_magnitude|pct1}} (interval {{exploratory.neg_p.median_site_pcross_magnitude:ci_low|pct1}} to {{exploratory.neg_p.median_site_pcross_magnitude:ci_high|pct1}}), similar to the increase case. PJM-adj's pooled bias turned negative, {{exploratory.neg_p.bias_pooled.pjmA.m0.1|pct1}} at a ten percent decrease and negative at {{exploratory.neg_p.sites_with_pjmA_bias_negative_at_m0.10|int}} of {{sample.n_sites|int}} sites, which is an underpayment of the participant. Expressed in energy, the median site break-even equals {{exploratory.x5.median_over_sites_kwh|.1f}} kWh over the shaped hours (site range {{exploratory.x5.min_site_kwh|.1f}} to {{exploratory.x5.max_site_kwh|.1f}} kWh).

\subsection{Secondary analyses: portfolios and screening}\label{sec:results:secondary}

Both analyses are descriptive and outside the Holm family. In the pre-specified portfolio analysis (H6), aggregating buildings lowered the pooled payment error ratio of PJM-adj from {{H6.pjmA.PER.n1|.2f}} for single buildings to {{H6.pjmA.PER.nbig|.2f}} for portfolios of twenty, with CG-adj similar (Fig.~\ref{figS:portfolio}). The CG-adj minus PJM-adj difference at twenty buildings was {{H6.contrast.level_big|.3f}} (interval {{H6.contrast.level_big:ci_low|.3f}} to {{H6.contrast.level_big:ci_high|.3f}}, {{H6.contrast.level_big:n_sites|int}} sites) and includes zero (Supplement Table~\ref{tabS:secondary}). The screening tree assigned {{screen.test.assigned_share.pjmA|pct1}} of the test buildings to PJM-adj and gave a pooled relMAE ratio of {{screen.test.pooled_relMAE_chosen_over_pjmA|.3f}} against PJM-adj, with {{screen.test.n_sites_better|int}} of {{sample.n_sites|int}} sites better. It gives no gain over PJM-adj in 2017 (null result in Supplement Table~\ref{tabS:screen}).
